% vim: ai sw=2 spell spelllang=cs

\begin{frame}{\textsc{Valgrind}}

  \begin{itemize}
    \item Sleduje (zejména) operace s pamětí (\emph{Memcheck})
    \begin{itemize}
      \item Úniky paměti
      \item Neinicializované a nezarovnané přístupy
      \item Použití paměti po \code{free}
      \item \dots
    \end{itemize}

    \pause

    \item Několik dalších modulů
    \begin{itemize}
      \item Callgrind: profiluje cache CPU
      \item Massif: profiluje použití haldy
      \item Další pro problémy se zámky atd.
    \end{itemize}

    \pause

    \item Spuštění: \cmd{valgrind --volby_valgrindu binarka --volby_binarky}
    \item Další parametry
    \begin{itemize}
      \item \cmd{\--tool}: výběr nástroje shora
      \item \cmd{\--leak-check}: sledovat úniky paměti
      \item \cmd{\-v}: podrobnější informace
    \end{itemize}

    \item Integrován do prostředí \qtc
  \end{itemize}

\end{frame}

\begin{frame}{Úkol}

  \heading{Práce s nástrojem \textsc{Valgrind}}
  \begin{enumerate}
    \iffimuni
      \item Projděte si \file{pb173-bin/09/memcheck.c}
    \fi
    \item Spusťte \file{memcheck} v nástroji \textsc{Valgrind}
    \begin{itemize}
      \item V \qtcu: Analyze $\rightarrow$ Valgrind Memory Analyzer
    \end{itemize}

    \item Řiďte se pokyny na konci výpisu
    \begin{itemize}
      \item Vypište si všechny podrobnosti
      \item Přidáváním navržených parametrů (,,Rerun with'')
      \item V \qtcu: Tools $\rightarrow$ Options $\rightarrow$ Analyzer
	$\rightarrow$ Show reachable and indirectly lost blocks
    \end{itemize}

    \item Opravte, aby si \textsc{Valgrind} nestěžoval
  \end{enumerate}

\end{frame}
